%% notes-tex/034-isobutanol-wetlab/sections/3-rnaseq.tex
%% SOURCE: hand-authored, resting on tables/t1-strain-panel.tex,
%% tables/t3-genotypes.tex and the two mirrored documents cited in
%% sections/1-panel.tex and sections/2-yko.tex.
%% The design fact that drives the selection, that Ll_ilvD in YZy452 sits behind
%% P_GAL10 and is therefore off in glucose, is quoted below.

\section{The First Sequencing Round\secstatus{todo}}
\label{sec:rnaseq}

The ask is a low-flux anchor, the best producer, and steps in between. The panel
supplies two independent versions of that, both for isobutanol, and one sentence
in the Methods is what turns the first of them into three points rather than
two.

\notesec{The promoter that makes a ladder} The dehydratase in YZy452 is not
constitutive.

\begin{quote}\small
Because the wild-type copy of \gene{Ll\_ilvD} in YZy452 is expressed from a
P$_{\text{GAL10}}$ promoter, growing the library in glucose ensures this
wild-type copy is repressed, and only the \gene{Ll\_ilvD} variants from the CEN
plasmid library (inserted downstream of a P$_{\text{TDH3}}$ promoter) are
expressed.
\end{quote}

\noindent
YZy452 also carries \gene{ilv3}$\Delta$, so the native dihydroxyacid dehydratase
is gone. In glucose, C0049 therefore has no dehydratase at all, C0050 supplies
the wild-type bacterial enzyme from a CEN plasmid, and C0051 supplies the
improved allele from the same backbone. Three levels of one step, one host, one
carbon source.

\notesec{Round one, five strains} Two matched contrasts, on two different steps
of the pathway: \gene{Ll\_ilvD} wild type against I433V at 3.1-fold, and
\gene{ILV6} wild type against V110E at 1.6-fold, with C0049 as the floor beneath
the first. \Cref{tab:roundone} gives the five, and what every tube needs if it is
picked up at all.

%% SOURCE: experiments/034-isobutanol-wetlab/scripts/strain_panel.py
%% Growth conditions from the Methods of zhangBiosensorBranchedchainAmino2022:
%% single colony -> 1 mL SC or SC-ura + 2% glucose overnight -> 10 uL into
%% 1 mL of the same in a 24-well plate -> 20 h -> spin, resuspend in 1 mL
%% of the same + 15% sugar -> seal -> 48 h, 30 C, 200 rpm -> OD600, HPLC.
\begin{table}[htbp]
\centering
\caption{\textbf{Round one, and how each tube is grown.} Medium is uracil dropout exactly when a CEN \gene{URA3} plasmid has to be held, and the fermentation carbon source is not a free choice for C0049. Every strain marked yes runs on the paper's ordinary 24-well protocol: single colony, shaker, plate reader, HPLC, with no sorting and no special illumination. Leucine must be present in all cases, because \gene{leu2}$\Delta$ makes these strains leucine auxotrophs; standard SC supplies it.}
\label{tab:roundone}
\footnotesize
\begin{tabular}{@{}lllll >{\raggedright\arraybackslash}p{50mm}@{}}
\toprule
ID & Strain & Medium & Fermentation & Pick and use & Round one \\
\midrule
C0044 & YZy313 & SC-ura & 15\% glucose & yes & yes -- \gene{ILV6} wild type, low arm of the 1.6-fold pair \\
C0045 & YZy314 & SC-ura & 15\% glucose & yes & yes -- ILV6V110E, high arm of the 1.6-fold pair \\
C0049 & YZy452 & SC & 15\% galactose & yes & yes -- ladder floor, dehydratase off on glucose \\
C0050 & YZy454 & SC-ura & 15\% glucose & yes & yes -- wild-type allele anchor \\
C0051 & YZy469 & SC-ura & 15\% glucose & yes & yes -- Ll\_ilvDI433V, 3.1-fold over C0050 \\
\midrule
C0043 & Yzy311 & SC-ura & 15\% glucose & yes & -- \\
C0046 & YZy148 & SC & 15\% glucose & yes & -- \\
C0047 & YZy148+LEU4WT & SC-ura & 15\% glucose & yes & -- \\
C0048 & YZy148+LEU4deltaS547 & SC-ura & 15\% glucose & yes & -- \\
C0052 & YZy91 & SC & 15\% glucose & yes & -- \\
C0053 & YZy502 & SC & 2\% glucose, dark & \textbf{no} & -- \\
\bottomrule
\end{tabular}
\end{table}


\notesec{No sorting anywhere in this} The protocol that produced every titer in
the paper is a single colony into 1 mL of SC or SC-ura with 2 percent glucose
overnight, 10 $\mu$L of that into 1 mL of the same in a 24-well plate, 20 h, then
a spin and a resuspension into 1 mL of the same medium with 15 percent sugar,
sealed, 48 h at 30 $^{\circ}$C and 200 rpm, then OD$_{600}$ on a plate reader and
HPLC on the supernatant. A shaker, a plate reader and an HPLC. Fluorescence-activated
cell sorting appears only in the path that isolates strains out of pooled
libraries, and those isolations are already done.

\notesec{Two conditions that are not free choices} C0049 ferments on galactose,
not glucose, for the P$_{\text{GAL10}}$ reason above, and on glucose it will read
as a dead strain rather than as a floor. C0053 is the one tube that is not pick
and use at all: its overnight needs constant blue light before production can run
in the dark, which is an illuminated-plate setup rather than a shaker.

Leucine has to be in the medium for all of them. The SI is direct about why:
"LEU2 deletion also causes leucine auxotrophy, requiring supplementation of this
amino acid in the growth medium". Standard SC supplies it, and a leucine dropout
would kill the panel.

\notesec{Measure titer alongside, do not carry it across} C0049's published
310 $\pm$ 15 mg/L is a galactose number. Whichever carbon source the round
settles on, the titer for each strain belongs in the same experiment as the
sequencing rather than taken from the paper, because the media here will not
match the media there well enough for the numbers to transfer.

\noindent
C0053 stays out of round one on the light grounds above, but it is worth
sequencing eventually for one reason: it is the direct parent of YZy505, the
paper's best isobutanol strain at 681 $\pm$ 29 mg/L, so a baseline taken now is
what makes that comparison possible later.

\notesec{What is not in the round} The isopentanol arm, C0046 through C0048, is
the cleanest three-point series in the box and C0048 holds the highest titer of
anything here at 842 $\pm$ 23 mg/L. It is the leucine branch. This program is
the valine branch, so those three are round two, and they move up only if a
dose-response in $\alpha$-isopropylmalate flux is wanted for its own sake.

C0043 YZy311 sits out because its nearest neighbor in the panel, C0044, differs
from it by both host and pathway, so a contrast between them would not isolate
anything. C0052 YZy91 sits out because it is the plasmid-free host of C0043
rather than a producer.

\notesec{Three requests worth making with the next shipment}

\begin{itemize}\setlength\itemsep{0.15em}
  \item \textbf{YZy121}, which Supplementary Table 1 gives as CEN.PK2-1C
    carrying the isobutanol-configured biosensor and nothing else. It is the
    reference every strain in this panel is missing: same lineage, same
    reporter, no deletions and no pathway.
  \item \textbf{YZy312}, Strain B, which completes the Fig.~3c 2$\times$2 that
    C0043, C0044 and C0045 are three quarters of. It is one transformation of
    C0052 with pYZ228.
  \item \textbf{YZy453}, YZy452 carrying an empty CEN URA3 plasmid. C0049
    carries no plasmid while C0050 and C0051 do, so the plasmid itself and its
    selection marker are currently confounded with the dehydratase.
\end{itemize}

\noindent
None of the three connects this panel to the knockout screen; for that reason
\cref{sec:yko} names yJM1 instead.
